\chapter{کارکرد‌ها}
\label{ch:functions}

\begin{goals}
\begin{itemize}
  \item کارکرد چیست و چرا برنامه را به کارکرد‌ها تقسیم می‌کنیم
  \item تعریف کارکرد و فراخوانی آن
  \item فرستادن مقدار به کارکرد با پارامتر
  \item گرفتن نتیجه از کارکرد با \fcode{بازگشت}
  \item مقدار پیش‌فرض پارامترها و محدوده‌ی متغیرها
\end{itemize}
\end{goals}

\section{یک نام برای یک کار}

از فصل اول با یک کارکرد کار می‌کنیم: \fcode{کارکرد آغازین}. حالا وقت آن است که
ببینیم «کارکرد» واقعاً چیست. \textbf{کارکرد} بلوکی از دستورهاست که یک نام
دارد. هر جا آن نام را بنویسیم (به این کار \textbf{فراخوانی} می‌گویند)،
دستورهای درون کارکرد اجرا می‌شوند. کارکرد آغازین هم همین است، با این تفاوت که
سلام آن را خودش، در آغاز برنامه، فراخوانی می‌کند.

فرض کنید در برنامه‌ای چند بار باید یک خط جداکننده چاپ کنیم. می‌توانیم هر
بار \fcode{چاپ "--------------------"} بنویسیم، یا این کار را یک بار در
کارکردی تعریف کنیم و هر جا لازم بود، نامش را صدا بزنیم:

\program{تعریف و فراخوانی کارکرد}
\begin{salamfa}
کارکرد خط جداکننده():
    چاپ "--------------------"
پایان

کارکرد آغازین:
    خط جداکننده()
    چاپ "فهرست خرید"
    خط جداکننده()
    چاپ "نان، پنیر، شیر"
    خط جداکننده()
پایان
\end{salamfa}

\begin{outputfa}
--------------------
فهرست خرید
--------------------
نان، پنیر، شیر
--------------------
\end{outputfa}

تعریف کارکرد با واژه‌ی \kwd{کارکرد} شروع می‌شود، بعد نام کارکرد می‌آید، بعد یک
جفت پرانتز و دونقطه، و در پایان \fcode{پایان}. نام کارکرد هم مثل نام متغیر
می‌تواند چند واژه‌ای باشد. فراخوانی، نوشتن نام کارکرد همراه با پرانتز است:
\fcode{خط جداکننده()}. پرانتزها لازم‌اند؛ بدون آن‌ها سلام فکر می‌کند از یک
متغیر حرف می‌زنید.

به دو نکته دقت کنید. اول این که کارکرد را \emph{بیرون} کارکرد آغازین تعریف
کردیم، در بالای برنامه؛ کارکرد‌ها همیشه در سطح بالای برنامه و کنار هم
تعریف می‌شوند، نه درون یکدیگر. دوم این که ترتیب تعریف مهم نیست؛ می‌توانستیم
\fcode{خط جداکننده} را زیر \fcode{کارکرد آغازین} هم بنویسیم. سلام اول همه‌ی
تعریف‌ها را می‌بیند و بعد از \fcode{کارکرد آغازین} شروع می‌کند.

\begin{note}[چرا کارکرد؟]
سه دلیل اصلی دارد. \textbf{تکرار نکردن}: کاری که در چند جا لازم است، یک بار
نوشته می‌شود و اگر روزی باید عوض شود، یک جا عوض می‌شود. \textbf{نام
گذاشتن}: خواندن \fcode{خط جداکننده()} بسیار آسان‌تر از خواندن
\fcode{چاپ "--------------------"} است؛ نام، مقصود را می‌گوید.
\textbf{تقسیم کار}: یک مسئله‌ی بزرگ را به چند کارکرد کوچک می‌شکنیم و هر
کدام را جداگانه می‌نویسیم و آزمایش می‌کنیم.
\end{note}

\section{پارامتر: فرستادن مقدار به کارکرد}

کارکرد \fcode{خط جداکننده} همیشه یک کار ثابت می‌کند. اما بیشتر کارکرد‌ها به
اطلاعاتی نیاز دارند تا کارشان را انجام دهند: نام کسی که باید به او سلام
کرد، عددی که باید مربعش را حساب کرد. این اطلاعات را از راه
\textbf{پارامتر} به کارکرد می‌فرستیم. پارامترها را درون پرانتز تعریف کارکرد
می‌نویسیم و برای هر یک باید \emph{ریختش} را هم بگوییم:

\program{کارکرد با پارامتر}
\begin{salamfa}
کارکرد سلام کن(نام: رشته):
    چاپ "سلام", نام + "! خوش آمدی."
پایان

کارکرد آغازین:
    سلام کن("سارا")
    سلام کن("رضا")
    دوست := "مینا"
    سلام کن(دوست)
پایان
\end{salamfa}

\begin{outputfa}
سلام سارا! خوش آمدی.
سلام رضا! خوش آمدی.
سلام مینا! خوش آمدی.
\end{outputfa}

\fcode{نام: رشته} می‌گوید این کارکرد یک پارامتر به نام \fcode{نام} از ریخت
رشته می‌خواهد. در هر فراخوانی، مقداری که درون پرانتز می‌نویسیم (به آن
\textbf{آرگومان} می‌گویند) در جعبه‌ی \fcode{نام} گذاشته می‌شود و کارکرد با
آن کار می‌کند. آرگومان می‌تواند یک مقدار مستقیم باشد، یک متغیر، یا حتی
یک محاسبه.

کارکرد می‌تواند چند پارامتر داشته باشد که با ویرگول جدا می‌شوند. هنگام
فراخوانی، آرگومان‌ها به همان ترتیب به پارامترها می‌رسند:

\program{چند پارامتر}
\begin{salamfa}
کارکرد ردیف علامت(تعداد: صحیح, علامت: رشته):
    گذرا خط := ""
    چرخه تعداد:
        خط += علامت
    پایان
    چاپ خط
پایان

کارکرد آغازین:
    ردیف علامت(۵, "*")
    ردیف علامت(۳, "#")
    ردیف علامت(۸, "=")
پایان
\end{salamfa}

\begin{outputfa}
*****
###
========
\end{outputfa}

\begin{warn}
ریخت آرگومان باید با ریخت پارامتر بخواند. فراخوانی \fcode{ردیف علامت("5", "*")}
خطاست، چون \code{"5"} رشته است و پارامتر اول عدد صحیح می‌خواهد. تعداد
آرگومان‌ها هم باید درست باشد؛ نه کمتر و نه بیشتر.
\end{warn}

\section{بازگشت: گرفتن نتیجه از کارکرد}

کارکرد‌های بالا کاری می‌کردند (چاپ)، اما چیزی به فراخواننده پس نمی‌دادند.
بسیاری از کارکرد‌ها باید نتیجه‌ای \emph{بدهند}: مساحت، قیمت نهایی، جواب یک
سؤال. برای این کار دو چیز لازم است: ریخت نتیجه را بعد از پرانتز اعلام
کنیم، و درون کارکرد با \kwd{بازگشت} نتیجه را پس بدهیم:

\program{کارکرد با مقدار بازگشتی}
\begin{salamfa}
کارکرد مساحت مستطیل(طول: صحیح, عرض: صحیح): صحیح:
    بازگشت طول * عرض
پایان

کارکرد آغازین:
    چاپ "مساحت اتاق:", مساحت مستطیل(۴, ۵)
    چاپ "مساحت حیاط:", مساحت مستطیل(۱۲, ۸)
    مجموع := مساحت مستطیل(۴, ۵) + مساحت مستطیل(۱۲, ۸)
    چاپ "مجموع:", مجموع
پایان
\end{salamfa}

\begin{outputfa}
مساحت اتاق: 20
مساحت حیاط: 96
مجموع: 116
\end{outputfa}

به سرآغاز کارکرد نگاه کنید: \fcode{کارکرد مساحت مستطیل(طول: صحیح, عرض: صحیح):
صحیح:}. بعد از پرانتز، یک دونقطه و ریخت نتیجه (\fcode{صحیح}) آمده و بعد
دونقطه‌ی همیشگی که بلوک را باز می‌کند. دستور \fcode{بازگشت طول * عرض}
مقدار را حساب می‌کند، آن را به فراخواننده می‌دهد و کارکرد همان‌جا پایان
می‌شود.

مهم‌ترین نکته این است که فراخوانی چنین کارکردی، خودش یک \emph{مقدار} است و
می‌توان آن را هر جا که یک مقدار می‌رود گذاشت: در \fcode{چاپ}، در یک
محاسبه، در تعریف یک متغیر، یا حتی به عنوان آرگومان کارکردی دیگر.

\subsection{کارکردی که بله یا نه می‌گوید}

کارکردی که نتیجه‌ی منطقی می‌دهد، در شرط‌ها بسیار خوش‌دست است:

\program{کارکرد منطقی}
\begin{salamfa}
کارکرد زوج است(عدد: صحیح): منطقی:
    بازگشت (عدد % ۲) == ۰
پایان

کارکرد آغازین:
    چرخه ۱ تا ۶ با عدد:
        اگر زوج است(عدد):
            چاپ عدد, "زوج"
        وگرنه:
            چاپ عدد, "فرد"
        پایان
    پایان
پایان
\end{salamfa}

\begin{outputfa}
1 فرد
2 زوج
3 فرد
4 زوج
5 فرد
6 زوج
\end{outputfa}

خط \fcode{اگر زوج است(عدد):} تقریباً مثل فارسی خوانده می‌شود. این یکی از
زیبایی‌های نام‌گذاری خوب کارکرد‌هاست. رسم خوبی است که نام کارکرد‌های منطقی را
به شکل پرسش بگذاریم: \fcode{زوج است}، \fcode{معتبر است}، \fcode{خالی است}.

\subsection{چند بازگشت در یک کارکرد}

کارکرد می‌تواند در چند جا \fcode{بازگشت} داشته باشد. به محض اجرای اولین
\fcode{بازگشت}، کارکرد تمام می‌شود و بقیه‌ی خط‌ها اجرا نمی‌شوند. این ویژگی
زنجیره‌های شرطی را کوتاه می‌کند:

\program{چند بازگشت}
\begin{salamfa}
کارکرد برچسب نمره(نمره: صحیح): رشته:
    اگر نمره >= ۱۸: بازگشت "عالی" پایان
    اگر نمره >= ۱۵: بازگشت "خوب" پایان
    اگر نمره >= ۱۰: بازگشت "قابل قبول" پایان
    بازگشت "مردود"
پایان

کارکرد آغازین:
    چاپ برچسب نمره(۱۹)
    چاپ برچسب نمره(۱۶)
    چاپ برچسب نمره(۱۲)
    چاپ برچسب نمره(۷)
پایان
\end{salamfa}

\begin{outputfa}
عالی
خوب
قابل قبول
مردود
\end{outputfa}

با نمره‌ی ۱۶، شرط اول نادرست است، شرط دوم درست است و کارکرد با «خوب»
برمی‌گردد؛ خط‌های بعدی هرگز دیده نمی‌شوند. مقایسه کنید با همان کار در فصل
\ref{ch:conditions} که با \fcode{وگرنه} نوشته بودیم.

\begin{warn}
سلام با کارکرد‌ها هم مثل متغیرها سخت‌گیر است: کارکردی که تعریف شده اما
هیچ‌جا فراخوانی نشده، خطاست. همچنین کارکردی که ریخت بازگشتی اعلام کرده،
باید در \emph{همه‌ی} مسیرها چیزی برگرداند. اگر \fcode{بازگشت "مردود"} آخر را حذف کنید، سلام خطا می‌دهد که
کارکرد ممکن است بدون بازگشت تمام شود. برعکس، کارکردی که ریخت بازگشتی ندارد
(مثل \fcode{سلام کن}) نباید مقداری برگرداند.
\end{warn}

\subsection{کارکرد با عدد اعشاری}

\program{تبدیل دما به شکل کارکرد}
\begin{salamfa}
کارکرد به فارنهایت(سلسیوس: اعشار۶۴): اعشار۶۴:
    بازگشت ((سلسیوس * ۹.۰) / ۵.۰) + ۳۲.۰
پایان

کارکرد آغازین:
    چاپ "آب یخ می‌زند:", به فارنهایت(۰.۰)
    چاپ "دمای بدن:", به فارنهایت(۳۷.۰)
    چاپ "آب می‌جوشد:", به فارنهایت(۱۰۰.۰)
پایان
\end{salamfa}

\begin{outputfa}
آب یخ می‌زند: 32
دمای بدن: 98.6
آب می‌جوشد: 212
\end{outputfa}

فرمولی که در فصل \ref{ch:loops} درون حلقه نوشته بودیم، حالا یک نام دارد
و هر جا لازم باشد قابل استفاده است. دقت کنید که آرگومان‌ها را با نقطه‌ی
اعشار نوشتیم (\code{۰.۰} نه \code{۰})، چون پارامتر اعشاری است.

\section{مقدار پیش‌فرض برای پارامتر}

گاهی یک پارامتر بیشتر وقت‌ها مقدار مشخصی دارد و فقط گاهی عوض می‌شود.
می‌توان در تعریف کارکرد برایش \textbf{مقدار پیش‌فرض} گذاشت؛ آن وقت
فراخواننده می‌تواند آن آرگومان را ننویسد:

\program{پارامتر با مقدار پیش‌فرض}
\begin{salamfa}
کارکرد سلام کن(نام: رشته, پیشوند: رشته = "سلام"):
    چاپ پیشوند, نام
پایان

کارکرد آغازین:
    سلام کن("سارا")
    سلام کن("دکتر اسماعیلی", "درود بر")
پایان
\end{salamfa}

\begin{outputfa}
سلام سارا
درود بر دکتر اسماعیلی
\end{outputfa}

پارامترهای پیش‌فرض‌دار باید بعد از پارامترهای معمولی بیایند.

\section{محدوده‌ی متغیرها}

متغیری که درون یک کارکرد تعریف می‌شود، فقط درون همان کارکرد وجود دارد. به
این ویژگی \textbf{محدوده} می‌گویند. کارکرد آغازین نمی‌تواند متغیرهای درون
\fcode{سلام کن} را ببیند و برعکس. پارامترها هم همین‌طورند: جعبه‌هایی که
با شروع کارکرد ساخته می‌شوند و با پایان آن از بین می‌روند.

نکته‌ی مهم‌تر این است که آرگومان‌ها به شکل \emph{کپی} به کارکرد می‌رسند. کارکرد
هر کاری با نسخه‌ی خودش بکند، متغیر فراخواننده دست‌نخورده می‌ماند:

\program{کارکرد کپی می‌گیرد}
\begin{salamfa}
کارکرد دو برابر کن(عدد: صحیح): صحیح:
    گذرا نتیجه := عدد
    نتیجه *= ۲
    بازگشت نتیجه
پایان

کارکرد آغازین:
    مقدار := ۲۱
    چاپ "نتیجه‌ی کارکرد:", دو برابر کن(مقدار)
    چاپ "مقدار اصلی:", مقدار
پایان
\end{salamfa}

\begin{outputfa}
نتیجه‌ی کارکرد: 42
مقدار اصلی: 21
\end{outputfa}

این همان چیزی است که در فصل \ref{ch:variables} درباره‌ی کپی شدن مقدارها
دیدیم. اگر کارکرد بخواهد نتیجه‌ای به بیرون بدهد، راه درستش \fcode{بازگشت}
است، نه دست‌کاری آرگومان.

\subsection{پایدارهای سراسری}

گاهی مقداری در چند کارکرد لازم است. \fcode{پایدار}‌هایی که در بالای برنامه و
بیرون همه‌ی کارکرد‌ها تعریف شوند، در همه جا در دسترس‌اند:

\program{پایدار سراسری}
\begin{salamfa}
پایدار مالیات := ۹

کارکرد قیمت نهایی(قیمت: صحیح): صحیح:
    بازگشت قیمت + ((قیمت * مالیات) / ۱۰۰)
پایان

کارکرد آغازین:
    چاپ قیمت نهایی(۱۰۰۰۰۰)
    چاپ قیمت نهایی(۲۵۰۰۰۰)
    چاپ "نرخ:", مالیات, "درصد"
پایان
\end{salamfa}

\begin{outputfa}
109000
272500
نرخ: 9 درصد
\end{outputfa}

بیرون از کارکرد‌ها فقط \fcode{پایدار} تعریف کنید، نه متغیر تغییرپذیر. متغیر
سراسری‌ای که هر کارکردی بتواند عوضش کند، پیدا کردن اشکال‌ها را بسیار سخت
می‌کند؛ بهتر است هر چه یک کارکرد لازم دارد را از راه پارامتر بگیرد و هر چه
تولید می‌کند را از راه بازگشت پس بدهد.

\begin{deep}[کارکردی که خودش را صدا می‌زند]
کارکرد می‌تواند درون خودش، خودش را فراخوانی کند. به این کار
\textbf{بازگشتی} می‌گویند و کارهای جالبی با آن می‌شود کرد. برنامه‌ی زیر
با همین ترفند شمارش معکوس می‌کند: هر فراخوانی، عددی یکی کمتر را
می‌فرستد، تا وقتی که به صفر برسد و زنجیره قطع شود. \fcode{بازگشت} بدون
مقدار، کارکرد بی‌نتیجه را همان‌جا تمام می‌کند. در این کتاب به این موضوع
بیشتر نمی‌پردازیم، اما بدانید که وجود دارد.
\end{deep}

\begin{salamfa}
کارکرد پس شماری(عدد: صحیح):
    اگر عدد == ۰:
        چاپ "پایان"
        بازگشت
    پایان
    چاپ عدد
    پس شماری(عدد - ۱)
پایان

کارکرد آغازین:
    پس شماری(۳)
پایان
\end{salamfa}
\begin{outputfa}
3
2
1
پایان
\end{outputfa}

\section{یک برنامه‌ی کامل‌تر}

صورت‌حساب یک رستوران، که هر بخش کارش را به یک کارکرد سپرده است:

\program{صورت‌حساب رستوران}
\begin{salamfa}
پایدار خدمات := ۱۰

کارکرد مبلغ سفارش(قیمت واحد: صحیح, تعداد: صحیح): صحیح:
    بازگشت قیمت واحد * تعداد
پایان

کارکرد اعمال تخفیف(مبلغ: صحیح, درصد: صحیح): صحیح:
    بازگشت مبلغ - ((مبلغ * درصد) / ۱۰۰)
پایان

کارکرد نمایش سطر(عنوان: رشته, مبلغ: صحیح):
    چاپ عنوان + ":", مبلغ, "تومان"
پایان

کارکرد آغازین:
    چلوکباب := مبلغ سفارش(۱۸۰۰۰۰, ۲)
    نوشابه := مبلغ سفارش(۱۵۰۰۰, ۳)
    جمع := چلوکباب + نوشابه
    حق خدمات := (جمع * خدمات) / ۱۰۰

    نمایش سطر("چلوکباب", چلوکباب)
    نمایش سطر("نوشابه", نوشابه)
    نمایش سطر("جمع", جمع)
    نمایش سطر("حق خدمات", حق خدمات)
    نمایش سطر("قابل پرداخت", اعمال تخفیف(جمع + حق خدمات, ۵))
پایان
\end{salamfa}

\begin{outputfa}
چلوکباب: 360000 تومان
نوشابه: 45000 تومان
جمع: 405000 تومان
حق خدمات: 40500 تومان
قابل پرداخت: 423225 تومان
\end{outputfa}

کارکرد آغازین این برنامه را بخوانید: تقریباً مثل یک متن فارسی است، چون هر
جزئیات محاسبه در کارکردی با نام گویا پنهان شده است. به آخرین خط دقت کنید:
نتیجه‌ی \fcode{اعمال تخفیف} مستقیم به عنوان آرگومان به \fcode{نمایش سطر}
داده شده است. کارکرد‌ها این‌گونه در هم می‌نشینند.

\begin{summary}
\begin{itemize}
  \item کارکرد بلوکی نام‌دار است: \fcode{کارکرد نام(پارامترها): ریخت بازگشتی:}
        \ldots\ \fcode{پایان}. فراخوانی با نام و پرانتز انجام می‌شود.
  \item کارکرد‌ها در سطح بالای برنامه تعریف می‌شوند؛ ترتیبشان مهم نیست.
  \item پارامترها با ریختشان نوشته می‌شوند: \fcode{نام: رشته}. آرگومان‌ها
        به ترتیب و به شکل کپی به آن‌ها می‌رسند.
  \item \fcode{بازگشت مقدار} نتیجه را پس می‌دهد و کارکرد را تمام می‌کند.
        فراخوانی چنین کارکردی خودش یک مقدار است.
  \item پارامتر می‌تواند مقدار پیش‌فرض داشته باشد.
  \item متغیرهای درون کارکرد فقط در همان کارکرد وجود دارند. برای مقدارهای
        مشترک، \fcode{پایدار} سراسری تعریف کنید.
\end{itemize}
\end{summary}

\section*{تمرین‌ها}

\exercise کارکردی به نام \fcode{بزرگتر} بنویسید که دو عدد صحیح بگیرد و
بزرگ‌ترشان را برگرداند. سپس با کمک آن، بزرگ‌ترین سه عدد را پیدا کنید.

\exercise کارکردی بنویسید که شعاع دایره (اعشاری) را بگیرد و مساحتش را
برگرداند (عدد پی را ۳٫۱۴۱۶ بگیرید). با یک حلقه، مساحت دایره‌هایی با شعاع
۱ تا ۵ را چاپ کنید.

\exercise کارکردی به نام \fcode{سال کبیسه است} بنویسید که یک سال شمسی
بگیرد و بگوید آیا کبیسه است یا نه. قاعده‌ی ساده‌شده: سالی کبیسه است که
باقی‌مانده‌ی تقسیمش بر ۳۳ یکی از ۱، ۵، ۹، ۱۳، ۱۷، ۲۲، ۲۶ یا ۳۰ باشد. سپس
برای سال‌های ۱۴۰۳ تا ۱۴۰۸ نتیجه را چاپ کنید.

\exercise کارکردی بنویسید که تعداد ثانیه بگیرد و آن را به شکل رشته‌ای مثل
\code{"1:02:05"} (ساعت:دقیقه:ثانیه) برگرداند. راهنمایی: از محاسبه‌ی فصل
\ref{ch:operators} و چسباندن رشته‌ها استفاده کنید.

\exercise برنامه‌ی جدول ضرب فصل \ref{ch:loops} را با کارکردی به نام
\fcode{سطر جدول} بازنویسی کنید که شماره‌ی سطر و تعداد ستون‌ها را بگیرد و
یک سطر جدول را چاپ کند. تعداد ستون‌ها مقدار پیش‌فرض ۱۰ داشته باشد.

\exercise کارکرد زیر یک اشکال دارد. بدون اجرا پیدایش کنید:
\begin{salamfa}
کارکرد قدرمطلق(عدد: صحیح): صحیح:
    اگر عدد < ۰:
        بازگشت -عدد
    وگرنه عدد > ۰:
        بازگشت عدد
    پایان
پایان
\end{salamfa}
